\documentclass[12pt]{article}
\usepackage[utf8]{inputenc}
\usepackage[T1]{fontenc}
\usepackage{hyperref}
\title{Talaud Language Resources for K-12 Curriculum Development}
\author{Compiled by Claude Sonnet 5 (claude-sonnet-5) for Andr\'as Kornai}
\date{}
\begin{document}
\maketitle

\section{Introduction}
Talaud (ISO 639-3 tld) is an Austronesian (Sangiric) language of the Talaud Islands, North Sulawesi, Indonesia. Wikipedia gives about 82,000 speakers and two dialects, Lami and Tirawata. Only limited material was found.

\section{Primary Educational Resources}
No primers, textbooks or school programmes were found.

\section{Linguistic and Reference Materials}
\begin{itemize}
\item Language overview: \url{https://en.wikipedia.org/wiki/Talaud_language}
\item Utsumi, ``A Dialectal Study of Talaud: Dialectal Classification and a Comparative Word List'' (Meisei University bulletin, 2016): \url{https://meisei.repo.nii.ac.jp/record/866/files/J-nbkNo024a05.pdf}
\item Austronesian Comparative Dictionary, Talaud entries: \url{https://acd.clld.org/languages/19277}
\item G. Bawole, \emph{Struktur Bahasa Talaud} (1981), a grammar in Indonesian listed in Wikipedia's further reading. No link was found.
\end{itemize}

\section{Religious and Community Resources}
No complete Talaud Bible was found. Global Recordings Network lists Bible lessons in the language: \url{https://globalrecordings.net/en/language/tld}. Joshua Project has a people-group profile: \url{https://joshuaproject.net/people_groups/15213}

\section{Digital and Audio Resources}
The Global Recordings Network page above lists audio recordings and a film in Talaud.

\section{Curriculum Development Context}
There is a word list and a dialect study but no school materials. A curriculum would need new material, built on the Bawole grammar and Utsumi's word list.
\end{document}